\documentclass[../../../include/open-logic-section]{subfiles}

\begin{document}
\olfileid{sth}{ord-arithmetic}{intro}

\olsection{引言}

在 \olref[ordinals][]{chap} 中，我们发展了序数理论。我们在 \olref[spine][]{chap} 中看到，可以将序数视为层级结构其余部分围绕其构建的中脊。但这并非序数唯一的角色。此外还有进行序数算术的任务。

我们早在 \olref[ordinals][idea]{sec} 讨论 $\omega$、$\omega+1$ 和 $\omega+\omega$ 时，就已经暗示过这一点。当时我们仅作了非形式的讨论；现在是时候严格地加以阐明了。不过应当指出，本章并没有太多哲学内容；只有技术性内容，外加一个（稍微）有趣的观察：我们可以用两种不同的方式做同一件事。

\end{document}